\documentclass[11pt,a4paper]{article}

\usepackage[utf8]{inputenc}
\usepackage[T1]{fontenc}
\usepackage[french]{babel}
\usepackage{geometry}
\usepackage{xcolor}
\usepackage{hyperref}
\usepackage{tabularx}
\usepackage{longtable}
\usepackage{enumitem}
\usepackage{titlesec}

\geometry{margin=2.4cm}
\hypersetup{
  colorlinks=true,
  linkcolor=blue!55!black,
  urlcolor=blue!55!black
}

\definecolor{chorusblue}{HTML}{1D4ED8}
\definecolor{chorusgray}{HTML}{F3F4F6}
\definecolor{chorusdark}{HTML}{111827}

\titleformat{\section}
  {\Large\bfseries\color{chorusblue}}
  {\thesection}{0.75em}{}

\titleformat{\subsection}
  {\large\bfseries\color{chorusdark}}
  {\thesubsection}{0.75em}{}

\setlist[itemize]{leftmargin=1.2cm}
\setlist[enumerate]{leftmargin=1.2cm}

\title{
  \vspace{-1.5cm}
  \textbf{\Huge Guide utilisateur}\\
  \vspace{0.3cm}
  \Large Chorus Pilot
}
\author{Projet Chorus}
\date{\today}

\begin{document}

\maketitle
\thispagestyle{empty}

\begin{center}
  \vspace{0.5cm}
  \begin{tabularx}{0.9\textwidth}{|X|}
    \hline
    \textbf{Objectif du document} \\
    Ce guide explique comment utiliser Chorus Pilot au quotidien pour
    consulter, préparer et envoyer des factures vers Chorus Pro. Il s'adresse à
    un utilisateur métier ou à un référent applicatif qui n'a pas besoin de
    connaître le code source ou les détails techniques internes. \\
    \hline
  \end{tabularx}
\end{center}

\newpage
\tableofcontents
\newpage

\section{Présentation rapide}

Chorus Pilot est une application qui aide à traiter les factures destinées
à Chorus Pro. Elle permet de consulter les factures, d'importer un document, de
vérifier les informations importantes et de transmettre la facture à Chorus.

L'application sert principalement à :

\begin{itemize}
  \item suivre les factures déjà connues ;
  \item préparer une nouvelle facture ;
  \item extraire les informations d'un PDF ;
  \item vérifier les montants, le fournisseur, le destinataire et le service ;
  \item envoyer la facture vers Chorus Pro ;
  \item consulter les statuts et les retours de traitement.
\end{itemize}

\section{À qui s'adresse ce guide ?}

Ce guide est destiné aux personnes qui utilisent l'application sans intervenir
sur son développement :

\begin{itemize}
  \item gestionnaires de factures ;
  \item agents administratifs ;
  \item référents Chorus ;
  \item utilisateurs chargés de contrôler ou soumettre les factures ;
  \item responsables souhaitant suivre l'activité.
\end{itemize}

Le guide évite volontairement le vocabulaire technique. Lorsqu'un terme
technique est nécessaire, il est expliqué simplement.

\section{Connexion à l'application}

\subsection{Ouvrir l'application}

Depuis votre navigateur, ouvrez l'adresse fournie par votre administrateur.
Selon l'installation, l'adresse peut ressembler à :

\begin{itemize}
  \item \texttt{http://localhost}
  \item \texttt{http://adresse-du-serveur}
  \item \texttt{http://chorusPilot.gigalis.local}
\end{itemize}

\subsection{Se connecter}

Sur l'écran de connexion :

\begin{enumerate}
  \item saisissez l'identifiant ou l'adresse e-mail fournie ;
  \item saisissez le mot de passe ;
  \item cliquez sur le bouton de connexion.
\end{enumerate}

Si la connexion échoue, vérifiez d'abord l'identifiant et le mot de passe. Si le
problème persiste, contactez la personne qui administre l'application.

\section{Navigation générale}

Après connexion, le menu latéral donne accès aux principales parties de
l'application.

\begin{longtable}{|p{0.3\textwidth}|p{0.62\textwidth}|}
  \hline
  \textbf{Menu} & \textbf{Utilisation} \\
  \hline
  Tableau de bord &
  Voir une synthèse de l'activité : nombre de factures, montants, statuts,
  répartition et dernières synchronisations. \\
  \hline
  Factures &
  Rechercher, filtrer et consulter les factures. \\
  \hline
  Nouvelle facture &
  Importer un document et préparer une soumission vers Chorus. \\
  \hline
  Clients publics &
  Consulter les destinataires publics connus. \\
  \hline
  Services destinataires &
  Rechercher les services associés aux clients publics. \\
  \hline
  Paramètres &
  Vérifier la configuration Chorus et tester la connexion. \\
  \hline
  Connexion / Profil &
  Accéder à la session utilisateur. \\
  \hline
\end{longtable}

\section{Tableau de bord}

Le tableau de bord donne une vue d'ensemble de l'activité.

Vous pouvez l'utiliser pour :

\begin{itemize}
  \item visualiser les volumes de factures ;
  \item suivre les montants concernés ;
  \item repérer les factures rejetées ou en attente ;
  \item filtrer par fournisseur, client public, statut ou période ;
  \item lancer une actualisation des données.
\end{itemize}

\subsection{Actualiser les données}

Le bouton d'actualisation permet de demander à l'application de récupérer les
données les plus récentes disponibles. Cette opération peut prendre quelques
instants selon le volume de factures.

Après actualisation, vérifiez la date de dernière synchronisation et les éventuels
messages d'erreur.

\section{Consulter les factures}

L'écran \textbf{Factures} permet de retrouver rapidement une facture.

\subsection{Rechercher une facture}

Vous pouvez rechercher une facture avec :

\begin{itemize}
  \item son numéro ;
  \item un nom de destinataire ;
  \item un statut ;
  \item une période de dépôt ;
  \item une période de facture ;
  \item un montant minimum ou maximum ;
  \item un numéro de bon de commande ;
  \item un mode de dépôt.
\end{itemize}

\subsection{Comprendre les statuts}

Les statuts indiquent l'état courant de la facture. Les libellés exacts peuvent
varier selon le retour Chorus, mais les grandes familles sont les suivantes :

\begin{longtable}{|p{0.3\textwidth}|p{0.62\textwidth}|}
  \hline
  \textbf{Type de statut} & \textbf{Signification} \\
  \hline
  En cours &
  La facture est connue et son traitement n'est pas terminé. \\
  \hline
  Déposée ou soumise &
  La facture a été transmise à Chorus. \\
  \hline
  Rejetée &
  La facture a été refusée ou n'a pas pu être traitée correctement. Une
  vérification est nécessaire. \\
  \hline
  Payée ou paiement &
  La facture est associée à une étape de paiement ou à un règlement. \\
  \hline
\end{longtable}

\subsection{Ouvrir le détail d'une facture}

Depuis la liste, ouvrez une facture pour consulter :

\begin{itemize}
  \item ses informations principales ;
  \item son destinataire ;
  \item ses montants ;
  \item son statut ;
  \item son historique ;
  \item les éventuels retours de Chorus.
\end{itemize}

\section{Créer une nouvelle facture}

L'écran \textbf{Nouvelle facture} sert à préparer une facture avant son envoi à
Chorus.

\subsection{Importer le document}

Commencez par déposer le fichier de facture. L'application accepte les documents
prévus par le projet, principalement :

\begin{itemize}
  \item les fichiers PDF ;
  \item les fichiers XML lorsque le flux le prévoit.
\end{itemize}

Après l'import, l'application tente de lire le document et d'extraire les
informations utiles.

\subsection{Vérifier les informations extraites}

Les informations détectées automatiquement doivent toujours être vérifiées.
Contrôlez en priorité :

\begin{itemize}
  \item le numéro de facture ;
  \item la date de facture ;
  \item le montant hors taxes ;
  \item le montant de TVA ;
  \item le montant toutes taxes comprises ;
  \item le fournisseur ;
  \item le client public destinataire ;
  \item le service destinataire ;
  \item les références de commande ou de marché si elles sont obligatoires.
\end{itemize}

Si une information est absente ou incorrecte, corrigez-la avant de poursuivre.

\subsection{Ajouter des pièces complémentaires}

Selon le cas, vous pouvez associer des documents complémentaires, par exemple :

\begin{itemize}
  \item bon de commande ;
  \item bon de livraison ;
  \item document annexe ;
  \item bordereau de suivi.
\end{itemize}

Chaque pièce doit être nommée clairement pour faciliter le contrôle ultérieur.

\subsection{Envoyer la facture}

Lorsque toutes les informations sont contrôlées :

\begin{enumerate}
  \item vérifiez une dernière fois les montants et le destinataire ;
  \item lancez la soumission ;
  \item attendez le retour de l'application ;
  \item consultez le message affiché après l'envoi.
\end{enumerate}

Si l'envoi réussit, l'application affiche un retour positif et peut indiquer une
référence Chorus. Si l'envoi échoue, un message d'erreur aide à identifier le
point à corriger.

\section{Utiliser les référentiels}

\subsection{Clients publics}

Le référentiel des clients publics permet de retrouver les structures
destinataires connues. Vous pouvez rechercher une structure par :

\begin{itemize}
  \item raison sociale ;
  \item code destinataire ;
  \item SIRET.
\end{itemize}

La consultation d'une fiche peut afficher des informations utiles sur la
structure, comme son identité ou son état d'activité.

\subsection{Services destinataires}

Le référentiel des services destinataires permet de trouver le service à utiliser
pour une facture. Ce point est important : une facture peut être rejetée si le
service destinataire est absent ou incorrect lorsque Chorus l'exige.

\subsection{Synchroniser l'annuaire}

La synchronisation de l'annuaire met à jour les clients publics et services
destinataires disponibles dans l'application.

À utiliser notamment :

\begin{itemize}
  \item lorsqu'un client public est introuvable ;
  \item lorsqu'un service destinataire semble manquant ;
  \item après une modification connue côté Chorus ;
  \item à la demande de l'administrateur applicatif.
\end{itemize}

\section{Paramètres}

L'écran \textbf{Paramètres} sert à vérifier que l'application est correctement
reliée à Chorus.

Un utilisateur non technique peut principalement contrôler :

\begin{itemize}
  \item l'environnement sélectionné : qualification ou production ;
  \item l'état du connecteur ;
  \item les valeurs métier utilisées par défaut ;
  \item le résultat du test de connexion.
\end{itemize}

\subsection{Qualification ou production}

\begin{itemize}
  \item \textbf{Qualification} : environnement de test, utilisé pour vérifier les
        flux sans impact réel ;
  \item \textbf{Production} : environnement réel, utilisé pour les factures à
        traiter officiellement.
\end{itemize}

Avant d'envoyer une facture réelle, assurez-vous que l'environnement attendu est
bien sélectionné.

\subsection{Tester la connexion}

Le test de connexion permet de vérifier que l'application peut communiquer avec
Chorus.

Si le test échoue, il ne faut pas insister sur les soumissions. Contactez la
personne responsable de la configuration avec le message affiché à l'écran.

\section{Que faire en cas de problème ?}

\subsection{La facture est rejetée}

Vérifiez les points suivants :

\begin{itemize}
  \item numéro de facture ;
  \item destinataire ;
  \item service destinataire ;
  \item montants HT, TVA et TTC ;
  \item références de commande ou de marché ;
  \item pièce jointe principale ;
  \item type de facture et mode de dépôt.
\end{itemize}

Après correction, suivez la procédure interne prévue pour une nouvelle soumission
ou une régularisation.

\subsection{Le PDF est mal lu}

Certains PDF sont plus difficiles à lire automatiquement, notamment lorsqu'ils
sont scannés comme des images ou générés avec une structure inhabituelle.

Dans ce cas :

\begin{itemize}
  \item contrôlez manuellement les champs extraits ;
  \item corrigez les montants et références ;
  \item vérifiez le destinataire et le service ;
  \item signalez le document si le problème se répète.
\end{itemize}

\subsection{Le destinataire est introuvable}

Essayez d'abord une recherche par SIRET ou par partie du nom. Si le résultat
reste introuvable :

\begin{itemize}
  \item lancez une synchronisation de l'annuaire si vous avez accès au bouton ;
  \item vérifiez que le SIRET est correct ;
  \item contactez le référent Chorus ou l'administrateur.
\end{itemize}

\subsection{La connexion à Chorus échoue}

Si le test de connexion échoue ou si les soumissions ne partent plus :

\begin{itemize}
  \item vérifiez que vous êtes dans le bon environnement ;
  \item notez le message d'erreur affiché ;
  \item évitez de multiplier les tentatives ;
  \item transmettez le message à l'administrateur.
\end{itemize}

\section{Bonnes pratiques}

\begin{itemize}
  \item Toujours vérifier les informations extraites automatiquement.
  \item Contrôler les montants avant soumission.
  \item Utiliser le bon environnement : qualification pour les tests, production
        pour les factures réelles.
  \item Rechercher le destinataire avec le SIRET lorsque le nom ne suffit pas.
  \item Ne pas ignorer un message d'erreur Chorus.
  \item Conserver les pièces justificatives utiles.
  \item Signaler les rejets récurrents pour améliorer le paramétrage.
\end{itemize}

\section{Glossaire simple}

\begin{longtable}{|p{0.28\textwidth}|p{0.64\textwidth}|}
  \hline
  \textbf{Terme} & \textbf{Définition} \\
  \hline
  Chorus Pro &
  Plateforme utilisée pour transmettre et suivre les factures destinées au secteur
  public. \\
  \hline
  PISTE &
  Accès technique permettant à l'application de communiquer avec Chorus Pro. \\
  \hline
  Fournisseur &
  Structure qui émet la facture. \\
  \hline
  Client public &
  Structure publique qui reçoit la facture. \\
  \hline
  Service destinataire &
  Service précis du client public auquel la facture doit être adressée. \\
  \hline
  SIRET &
  Identifiant officiel d'un établissement. \\
  \hline
  Payload &
  Ensemble des informations envoyées à Chorus. Pour l'utilisateur, cela correspond
  simplement au dossier de données préparé par l'application. \\
  \hline
  Synchronisation &
  Action de mettre à jour les données locales avec les données connues côté Chorus. \\
  \hline
  Rejet &
  Refus ou échec de traitement d'une facture. Une correction est généralement
  nécessaire. \\
  \hline
\end{longtable}

\section{Conclusion}

Chorus Pilot accompagne l'utilisateur dans les étapes importantes du
traitement d'une facture : recherche, préparation, contrôle, soumission et suivi.
L'application automatise une partie du travail, mais la vérification humaine reste
essentielle avant chaque envoi à Chorus Pro.

\end{document}
